\section{深度学习的局限性}

\subsection{重新审视泛化：随机标签实验}

Amini 介绍了一个经典的实验——来自 Zhang 等人 2017 年发表在 ICLR 的论文 \textit{Understanding Deep Neural Networks Requires Rethinking Generalization}。

\begin{knowledgebox}{实验设计}
研究者取 ImageNet 数据集中的图片及其标签，然后对每张图片独立地掷一个 $k$ 面骰子（$k$ 为类别总数），将原有标签替换为随机结果。这样做的后果是：同一类别的两张图可能被分配到完全不同的标签，标签与图像内容之间失去了所有语义关联。
\end{knowledgebox}

\begin{figure}[H]
\centering
\includegraphics[width=0.85\textwidth]{slides-images/slide-19.jpg}
\caption{随机标签实验：图像标签被完全打乱\protect\footnotemark}
\end{figure}
\footnotetext{Slides 第19页。引用 Zhang+ ICLR 2017。}

\subsubsection{实验结果}

研究者在不同程度的标签随机化下训练深度神经网络，观察训练集和测试集的表现：

\begin{figure}[H]
\centering
\includegraphics[width=0.85\textwidth]{slides-images/slide-22.jpg}
\caption{关键发现：即使标签完全随机，训练精度仍可达到 100\%\protect\footnotemark}
\end{figure}
\footnotetext{Slides 第22页。红色虚线标注训练精度始终接近 100\%。}

\begin{importantbox}{核心发现：深度网络的惊人拟合能力}
\begin{itemize}
    \item \textbf{测试精度}：随着标签随机化程度增加，测试集精度持续下降（符合直觉）。
    \item \textbf{训练精度}：无论标签多随机，深度网络始终可以在训练集上达到接近 100\% 的准确率。
\end{itemize}
这意味着现代深度网络拥有足够的容量（capacity）来``记住''整个训练集，即使数据中完全不存在可学习的模式。
\end{importantbox}

\subsection{函数逼近器的双刃剑}

基于上述实验，Amini 进一步用函数拟合的视角来阐释这一现象：

\begin{figure}[H]
\centering
\includegraphics[width=0.85\textwidth]{slides-images/slide-25.jpg}
\caption{神经网络作为函数逼近器：在有训练数据的区域表现出色\protect\footnotemark}
\end{figure}
\footnotetext{Slides 第25页。}

神经网络可以非常好地拟合训练数据点附近的函数值。给定一个新的数据点（紫色），如果它落在训练数据的分布范围内，网络可以给出合理的预测。但关键问题是：

\begin{figure}[H]
\centering
\includegraphics[width=0.85\textwidth]{slides-images/slide-28.jpg}
\caption{数据分布外的行为不可预测\protect\footnotemark}
\end{figure}
\footnotetext{Slides 第28页。}

\begin{warningbox}{Out-of-Distribution 问题}
当输入数据超出训练分布（out-of-distribution, OOD）时，神经网络的行为是不可预测的。在没有训练数据的区域，拟合函数可能产生任意的输出值——而网络不会``告诉你''它不确定。这就引出了一个核心问题：\textbf{我们如何知道模型不知道什么？}
\end{warningbox}

\subsection{深度学习不是炼金术}

\begin{figure}[H]
\centering
\includegraphics[width=0.85\textwidth]{slides-images/slide-29.jpg}
\caption{Deep Learning $\neq$ Alchemy\protect\footnotemark}
\end{figure}
\footnotetext{Slides 第29页。引用 U. Muller, 6.S191 2018。}

Amini 强调：深度学习\textbf{不是}一个可以把任何东西丢进去就能得到完美输出的``魔法黑箱''。经典的 ``garbage in, garbage out''（垃圾进，垃圾出）原则完全适用。部署 AI 解决方案之前需要考虑两个核心问题：
\begin{enumerate}
    \item 这个任务是否真的适合用深度学习来解决？
    \item 你能为这个任务收集和整理出高质量的数据吗？
\end{enumerate}

\subsection{失败模式一：数据偏差}

\begin{figure}[H]
\centering
\begin{subfigure}[b]{0.48\textwidth}
    \includegraphics[width=\textwidth]{slides-images/slide-30.jpg}
    \caption{训练图像中的狗大量带有伸舌头的姿态}
\end{subfigure}
\hfill
\begin{subfigure}[b]{0.48\textwidth}
    \includegraphics[width=\textwidth]{slides-images/slide-31.jpg}
    \caption{CNN 为黑白图片着色时出现的粉色区域}
\end{subfigure}
\caption{数据偏差导致的失败：图像着色模型\protect\footnotemark}
\end{figure}
\footnotetext{Slides 第30--31页。引用 P. Isola 6.869。}

课程举了一个生动的例子：训练一个 CNN 将黑白照片转换为彩色照片。当输出的狗照片下巴区域出现了粉色斑块时，原因是训练集中大量狗的照片都是伸着粉色舌头的姿态，模型学到了``狗的下巴附近应该是粉色的''这一统计关联——而非真正理解了颜色的含义。

\subsection{失败模式二：安全关键场景的不确定性}

\begin{figure}[H]
\centering
\includegraphics[width=0.85\textwidth]{slides-images/slide-32.jpg}
\caption{Tesla Autopilot 致命事故案例\protect\footnotemark}
\end{figure}
\footnotetext{Slides 第32页。引用 ABC News。}

Amini 提到了一起真实的自动驾驶致命事故：一辆 Tesla 在同一路段多次出现向隔离带偏转的现象，最终导致了碰撞。调查发现，该路段的隔离带是在训练数据采集之后才建造的——自动驾驶系统遇到了一个\textbf{训练分布之外}的场景。

\begin{figure}[H]
\centering
\includegraphics[width=0.85\textwidth]{slides-images/slide-33.jpg}
\caption{不确定性在安全关键应用中至关重要\protect\footnotemark}
\end{figure}
\footnotetext{Slides 第33页。}

\begin{importantbox}{深度学习中的不确定性}
在以下场景中，可靠地检测和量化不确定性尤为重要：
\begin{itemize}
    \item \textbf{安全关键应用}：自动驾驶、医疗诊断、人脸识别
    \item \textbf{数据质量问题}：类别不平衡（imbalance）、数据噪声（data noise）
    \item \textbf{分布外检测}：当输入数据与训练分布显著不同时，模型应当``知道自己不知道''
\end{itemize}
\end{importantbox}

\subsection{失败模式三：对抗攻击}

对抗攻击（Adversarial Attacks）是深度学习中一个非常经典且重要的安全问题。

\begin{figure}[H]
\centering
\includegraphics[width=0.85\textwidth]{slides-images/slide-35.jpg}
\caption{对抗攻击示例：Temple (97\%) $\rightarrow$ Ostrich (98\%)\protect\footnotemark}
\end{figure}
\footnotetext{Slides 第35页。}

\subsubsection{对抗攻击的数学原理}

回顾梯度下降的核心公式：

$$
W \leftarrow W - \eta \frac{\partial J(W, x, y)}{\partial W}
$$

其中 $W$ 是权重，$\eta$ 是学习率，$J$ 是损失函数，$(x, y)$ 是固定的输入和标签。训练时，我们\textbf{固定输入 $x$ 和标签 $y$}，优化权重 $W$ 来最小化损失。

\begin{figure}[H]
\centering
\includegraphics[width=0.85\textwidth]{slides-images/slide-40.jpg}
\caption{对抗图像的生成公式\protect\footnotemark}
\end{figure}
\footnotetext{Slides 第40页。}

对抗攻击则完全翻转了这个过程：

$$
x \leftarrow x + \eta \frac{\partial J(W, x, y)}{\partial x}
$$

\begin{importantbox}{对抗攻击 vs. 正常训练}
\begin{itemize}
    \item \textbf{正常训练}：固定输入 $(x, y)$，优化权重 $W$ 来\textbf{最小化}损失
    \item \textbf{对抗攻击}：固定权重 $W$ 和标签 $y$，修改输入 $x$ 来\textbf{最大化}损失
\end{itemize}
关键区别：梯度的计算对象从 $\partial W$ 变成了 $\partial x$，优化方向从减号变成了加号。
\end{importantbox}

\subsubsection{物理世界中的对抗样本}

对抗攻击不仅存在于数字空间，还可以在物理世界中实现。Amini 介绍了 MIT 学生利用 3D 打印技术制造的物理对抗样本——一只经过精心设计纹理的 3D 打印乌龟，在各种角度和光照条件下，都被分类器识别为``步枪''而非``乌龟''。

\begin{figure}[H]
\centering
\includegraphics[width=0.85\textwidth]{slides-images/slide-42.jpg}
\caption{3D 打印的物理对抗样本：乌龟被识别为步枪\protect\footnotemark}
\end{figure}
\footnotetext{Slides 第42页。引用 Athalye+ ICML 2018。}

\begin{warningbox}{对抗攻击的现实威胁}
对抗样本的存在意味着：在安全关键场景（如自动驾驶的交通标志识别、安防系统的人脸识别），精心设计的微小扰动就可能导致系统做出完全错误的判断。这不仅是学术问题，更是实际部署 AI 系统时必须考虑的安全风险。
\end{warningbox}

\subsection{局限性总结}

\begin{figure}[H]
\centering
\includegraphics[width=0.85\textwidth]{slides-images/slide-46.jpg}
\caption{神经网络局限性完整列表\protect\footnotemark}
\end{figure}
\footnotetext{Slides 第46页。}

课程总结了神经网络的主要局限性：

\begin{table}[H]
\centering
\begin{tabular}{ll}
\toprule
\textbf{局限性} & \textbf{说明} \\
\midrule
数据饥渴 (Data hungry) & 通常需要百万级训练样本 \\
计算密集 (Computationally intensive) & 训练和部署都需要 GPU \\
对抗脆弱 (Adversarial examples) & 容易被精心设计的扰动欺骗 \\
算法偏见 (Algorithmic bias) & 可能放大数据中的社会偏见 \\
不确定性表征差 (Uncertainty) & 难以知道模型何时``不知道'' \\
黑箱不可解释 (Black boxes) & 难以理解和信任 \\
需要专家知识 (Expert knowledge) & 架构设计和调参依赖经验 \\
\textbf{难以编码结构} (Encode structure) & 难以在学习中融入先验知识 \\
\textbf{外推困难} (Extrapolation) & 难以超越训练数据的分布范围 \\
\bottomrule
\end{tabular}
\caption{神经网络的主要局限性}
\end{table}

\begin{knowledgebox}{局限性孕育新前沿}
Amini 指出：作为技术人员和科学家，局限性恰恰代表着发展机遇。课程后半部分介绍的 Diffusion Models 和 LLMs 正是对``外推困难''和``编码结构困难''这两个核心挑战的回应——通过建模数据分布本身，基础模型（Foundation Models）有望实现更强大的泛化能力。
\end{knowledgebox}

\subsection{本章小结}

深度学习的局限性涵盖多个层面：泛化能力不足（随机标签实验揭示了惊人的过拟合能力）、数据偏差导致的系统性错误、安全关键场景中的不确定性挑战，以及对抗攻击对模型鲁棒性的威胁。理解这些局限性是负责任地使用 AI 技术的前提，也是推动下一代算法发展的动力。

%% =====================================================
